\documentclass[11pt]{article}

\usepackage[margin=1in]{geometry}
\usepackage{amsmath,amssymb,amsthm}
\usepackage{booktabs}
\usepackage{microtype}
\usepackage{natbib}
\usepackage[hidelinks]{hyperref}

\newtheorem{proposition}{Proposition}
\newcommand{\E}{\mathbb{E}}
\newcommand{\R}{\mathbb{R}}

\title{Residual Surrogates for Nonlinear DSGE Estimation:\\
An Auditable Target-Corrected Workflow}
\author{Anonymous working draft}
\date{October 23, 2026}

\begin{document}
\maketitle

\begin{abstract}
Repeated nonlinear model solution can dominate likelihood-based estimation of
dynamic stochastic general equilibrium models. We develop an auditable workflow
that learns the discrepancy between a first-order reduced-order model and
nonlinear extended-path solutions, evaluates the resulting profile inversion objective
in JAX, and uses a cheap differentiable density to construct Hamiltonian
proposals. A Metropolis test at each trajectory endpoint can preserve any
declared target density that is evaluated there. This correction is established
MCMC theory; our contribution is its integration with residual learning,
parameter-dependent DSGE solution, cross-implementation checks, and a frozen
evidence protocol. A preserved engineering pilot exercises the complete
23-equation nonlinear Gal\'i model while estimating two policy parameters. It
accepts 80 order-zero extended-path targets with maximum equation residual
$9.92\times10^{-9}$, reduces held-out state error relative to the first-order
model by 53.9 percent on average, and matches Python and JAX evaluations of the
same surrogate objective within $1.71\times10^{-13}$. These results establish
pipeline execution, not posterior fidelity or speedup: the pilot uses
deterministic extended path, and its corrected HMC target remains a surrogate
objective. We therefore state the confirmatory tests required before the method
can support substantive inference claims.
\end{abstract}

\section{Introduction}

Likelihood-based estimation of nonlinear DSGE models repeats a difficult
numerical task at many parameter values. Each likelihood evaluation can require
a steady state, a local solution, nonlinear conditional trajectories, shock
inversion, and measurement-error evaluation. The same work then appears inside
optimization or posterior simulation. A neural surrogate is attractive only if
it reduces this repeated cost without hiding the numerical failures that matter
for inference.

This paper develops a workflow around that requirement. The method first solves
the complete nonlinear equation system at a designed set of parameter, state,
and shock inputs. It then fits a residual network to the gap between those
solutions and a first-order reduced-order model (ROM). The ROM supplies global
economic structure and a strong baseline; the network learns only the remaining
nonlinearity. During posterior simulation, a cheap differentiable density can
generate Hamiltonian trajectories, while a separately declared target density
is evaluated at the endpoints. When the proposal map is reversible and volume
preserving, the endpoint Metropolis test preserves that declared target.

The central contribution is the integrated estimation and validation design,
not the generic observation that an approximate Hamiltonian may be corrected by
an exact Metropolis test. That result is standard \citep{neal2011hmc}, and
surrogate-gradient HMC has been studied with both random-basis and neural
surrogates \citep{zhang2017surrogate,li2019nnghmc}. The contribution must instead
be judged by four empirical questions. Does the residual surrogate improve on
the ROM outside its training parameter values? Does it preserve likelihood and
posterior geometry in the inference-relevant region? Does endpoint correction
target the density claimed in the paper? And does the complete workflow improve
effective samples per second after training and compilation costs are counted?

We answer only the first, engineering-level version of these questions in the
preserved pilot. The full confirmatory exercise is deliberately separated from
the pilot. This separation prevents a successful software run from being
mistaken for publication evidence.

\section{Model, solution objects, and likelihoods}

Let $\theta\in\Theta$ denote structural parameters, $s_t$ predetermined states,
$x_t$ endogenous variables, and $\varepsilon_t$ structural shocks. Write the
nonlinear equilibrium system as
\begin{equation}
  F(x_{t-1},x_t,\E_t x_{t+1},\varepsilon_t;\theta)=0.
  \label{eq:equilibrium}
\end{equation}
For a given parameter vector and shock design, a nonlinear extended-path solver
produces a trajectory $x^{F}_{1:T}(\theta)$ from the full-order model (FOM). A
first-order solution produces the cheaper trajectory $x^{R}_{1:T}(\theta)$. We
fit a network $r_{\psi}$ to the residual
\begin{equation}
  \Delta x_t(\theta)=x^{F}_t(\theta)-x^{R}_t(\theta),
  \qquad
  \widehat{x}^{F}_t(\theta)=x^{R}_t(\theta)+r_{\psi}(z_t,\theta),
  \label{eq:residual-surrogate}
\end{equation}
where $z_t$ contains the states, shocks, and other features available at the
prediction point. Parameter vectors, rather than individual rows, are assigned
to training or validation. This prevents trajectories from the same parameter
value from appearing on both sides of the split.

The target nonlinear posterior for data $y_{1:T}$ is
\begin{equation}
  \pi_F(\theta\mid y)\propto p(\theta)L_F(y_{1:T}\mid\theta),
  \label{eq:exact-posterior}
\end{equation}
where $L_F$ denotes an independently validated nonlinear likelihood with a
declared treatment of latent shocks and every required change-of-variables
term. The present implementation instead defines a surrogate inversion objective
$\widetilde L_{\psi}$ by replacing nonlinear trajectories with
\eqref{eq:residual-surrogate}, optimizing latent shocks, and evaluating the
shock and measurement penalties at that optimizer. Unless an integration,
Laplace adjustment, or required change-of-variables term is supplied, this is a
profile or pseudo-likelihood objective rather than a marginal likelihood. The
two objects are therefore scientifically distinct. Agreement between Python and
JAX evaluations of $\widetilde L_{\psi}$ is implementation parity; it is not
evidence that $\widetilde L_{\psi}=L_F$.

More explicitly, let $h_{\psi}(s_t,\varepsilon_t,\theta)$ map a state and
structural shock into the observables implied by the ROM plus residual
surrogate. At each date, the implementation computes
\begin{equation}
  \widehat\varepsilon_t(\theta)
  =\arg\min_{\varepsilon}
  \left\{
  \lVert R^{-1/2}[y_t-h_{\psi}(s_t,\varepsilon,\theta)]\rVert_2^2
  +\lVert Q^{-1/2}\varepsilon\rVert_2^2
  \right\},
  \label{eq:shock-profile}
\end{equation}
using a damped Gauss--Newton iteration, and then sums the Gaussian penalties
evaluated at $\widehat\varepsilon_t(\theta)$. Here $R$ and $Q$ are diagonal in
the pilot. The state passed to date $t+1$ follows the configured ROM or
surrogate replay. Equation~\eqref{eq:shock-profile} makes the missing operation
clear: the latent shocks are optimized, not integrated out. It also clarifies
why agreement between two implementations of this calculation is a software
check rather than likelihood validation.

The benchmark model is the complete nonlinear New Keynesian system distributed
with the MacroModelling.jl examples and associated with the model in
\citet{gali2015}. It contains 23 equations, 16 declared parameters, and three
structural shocks. The reduced pilot estimates the Taylor-rule responses to
inflation and output, $\phi_\pi$ and $\phi_y$, while retaining every model
equation. Reducing the parameter dimension makes a numerical posterior reference
feasible; it does not turn the model into the familiar three-equation
linearization.

\subsection{Offline and online computations}

The offline stage draws complete parameter vectors, solves a parameter-specific
steady state and first-order ROM, and generates nonlinear extended-path targets
conditional on designed state and shock paths. Each attempted target retains a
case identifier, solver status, residual norm, and parameter-group identifier.
Only the training partition determines feature scaling and network weights.
Validation partitions are disjoint in complete parameter vectors, not merely in
individual rows. The residual network and every preprocessing constant are
then frozen and serialized with the model identity and source hashes.

The online stage transforms an unconstrained vector $q$ to bounded structural
parameters, adds the transform Jacobian to the declared log target, rebuilds
the parameter-dependent objects required by the selected runtime, and evaluates
either the profile objective or a nonlinear reference. A surrogate-gradient
trajectory costs many cheap force evaluations but only one new target
evaluation per transition. This accounting---rather than draws per second
alone---is central when the endpoint target contains a nonlinear solution.

Extended path traces back to \citet{fair1983extended}. Deterministic extended
path sets future innovations to their conditional means and therefore does not
capture Jensen effects in nonlinear expectations. Stochastic extended path
approximates those conditional expectations over future shocks
\citep{adjemian2025sep}. This is why positive branching order, multiple nodes,
and active variation in every shock are hard requirements for confirmatory
evidence, even though order zero is useful for debugging deterministic
nonlinearity.

\section{Target-corrected surrogate trajectories}

Let $q$ be an unconstrained parameterization of $\theta$, let $p$ be Gaussian
momentum with kinetic energy $K(p)$, and let $\pi_T(q)$ be the density the
sampler is intended to preserve. A cheaper differentiable density $\pi_P(q)$
defines the force used in $L$ leapfrog steps. Denote the resulting deterministic
map by $(q',p')=\Phi_{\pi_P}(q,p)$. The acceptance probability is
\begin{equation}
  \alpha\{(q,p),(q',p')\}
  =\min\left\{1,
  \exp\left[-H_T(q',p')+H_T(q,p)\right]\right\},
  \quad
  H_T(q,p)=-\log\pi_T(q)+K(p).
  \label{eq:acceptance}
\end{equation}

\begin{proposition}[Target invariance]
Suppose $\Phi_{\pi_P}$ is reversible after momentum negation and volume
preserving, and $\pi_T$ is evaluated at both endpoints used in
\eqref{eq:acceptance}. Then the Metropolis transition leaves $\pi_T$ invariant,
even when $\pi_P\neq\pi_T$.
\end{proposition}

\noindent
The proposition follows from the standard deterministic-proposal Metropolis
argument for leapfrog HMC \citep{neal2011hmc}. It places two practical
restrictions on the implementation. First, the surrogate and every preprocessing
transformation used in the force must be frozen during production sampling.
Second, ``exact'' always modifies the declared target. If
$\pi_T\propto p(\theta)\widetilde L_{\psi}$, the chain is exact for the
surrogate-defined posterior, not for \eqref{eq:exact-posterior}. Exact nonlinear
posterior claims require endpoint evaluation of $L_F$ itself.

For bounded parameters $\theta_j\in(a_j,b_j)$, the implementation samples
$q_j\in\mathbb{R}$ and uses
\begin{equation}
  \theta_j=a_j+(b_j-a_j)\operatorname{logit}^{-1}(q_j),
  \qquad
  \log J(q)=\sum_j\left[
    \log(b_j-a_j)+\log\sigma(q_j)+\log\{1-\sigma(q_j)\}
  \right].
  \label{eq:bounded-transform}
\end{equation}
Thus the endpoint target is defined in $q$ coordinates and must include
$\log J(q)$. Invalid steady states, indeterminate local solutions, failed
nonlinear paths, or non-finite objective values must be rejected as points
outside support. A large finite penalty is not equivalent: it creates an
artificial plateau with positive density. The preserved pilot still uses such a
sentinel and therefore cannot support an exact structural-posterior claim.

The sampler has three implementation-level regression checks. When proposal and
target densities coincide, its draws and accept/reject decisions match the
ordinary static-HMC implementation under the same key. Under a deliberately
shifted and rescaled proposal force, a fixed-seed multi-chain test recovers the
known moments of a standard-normal target after endpoint correction. Finally,
non-finite candidate targets receive zero acceptance probability. These tests
check kernel plumbing; they do not replace convergence and posterior-fidelity
evidence for the DSGE application.

This distinction is especially relevant for differentiable state-space and
DSGE estimation, where automatic differentiation improves computation but does
not resolve approximation error in the target \citep{childers2022differentiable}.
Related work has made nonlinear DSGE likelihoods more tractable
\citep{kollmann2017tractable}; our evaluation therefore treats accuracy against
an independently computed nonlinear reference as a primary outcome.

\section{Validation design}

The confirmatory design has two tiers. Tier 1 estimates $(\phi_\pi,\phi_y)$ and
uses a dense posterior grid. Grid resolution is doubled until posterior moments
and interval endpoints stabilize. This tier isolates numerical correctness from
high-dimensional MCMC behavior. Tier 2 estimates a prespecified four-to-six
parameter subset and measures scaling with four independently initialized
chains. The paper will report Tier 2 only after Tier 1 passes.

Training, validation, and posterior-region designs are disjoint in parameter
vectors, shock sequences, and trajectories. Validation includes regular
prior-interior points, posterior-region points, and a boundary-stress stratum.
The last stratum is reported separately rather than deleted through a stability-
margin screen. For stochastic extended path, evidence runs use positive
branching order, multiple quadrature nodes, and nonzero coverage of every
declared shock. Common case identifiers are evaluated by both the Python and
MacroModelling.jl implementations.

The primary accuracy outcomes are scaled equilibrium residuals; held-out,
per-output trajectory errors; paired improvement over the ROM with uncertainty;
centered log-likelihood error; rank reversals near the mode; and normalized
gradient error away from zero-gradient cases. Posterior comparison reports
differences relative to pooled Monte Carlo error, interval endpoint differences,
and a distributional diagnostic. Sampling reports rank-normalized split
$\widehat R$, bulk and tail effective sample sizes (ESS), ESS per gradient
evaluation, and ESS per second.

Timing uses the same workload, precision, and hardware for both methods.
Compilation, host callbacks, device transfers, and training are reported
separately. Warm likelihood microbenchmarks use repeated synchronized calls;
end-to-end comparisons use independent processes. The current Schur QME path
invokes a sequential SciPy host callback, so a GPU label does not imply a
GPU-native likelihood. The available GPU-native matrix-sign alternative must
pass Gal\'i-specific solution and likelihood parity before it enters a speed
claim.

The protocol is frozen before confirmatory execution. It records the repository
revision, dirty-worktree state, environment, command, seeds, hardware, precision,
and content hashes for every evidence artifact. Numerical failures remain in
attempted-case denominators and are stratified by prior-interior,
posterior-region, and boundary-stress designs. Machine-readable evidence is the
source of manuscript tables; hand-entered numerical cells are not permitted.

\section{Preserved engineering pilot}

% Generated by benchmarks/build_methods_paper_evidence.py. Do not edit.
\begin{table}[t]
\centering
\caption{Preserved full-model Galí engineering pilots}
\label{tab:pilot-results}
\scriptsize
\setlength{\tabcolsep}{3pt}
\begin{tabular}{llrrrrrr}
\toprule
Run & Target runtime & $N_\theta \times T$ & Max. residual & Val. RMSE & Corrected & Accept. & Draws/s \\
\midrule
Quick smoke & fixed-reference & 4 $\times$ 4 & 7.21e-09 & 9.17e-05 & No & 1.000 & 0.378 \\
Corrected resource & full-jax & 4 $\times$ 4 & 7.21e-09 & 9.17e-05 & Yes & 1.000 & 0.0267 \\
Default pilot & fixed-reference & 8 $\times$ 10 & 9.92e-09 & 7.86e-06 & No & 1.000 & 8 \\
\bottomrule
\end{tabular}

\vspace{0.4em}
\parbox{0.98\textwidth}{\footnotesize\textit{Notes:} All runs use the complete 23-equation nonlinear Galí source but estimate only two policy-rule parameters. $N_\theta$ is the number of parameter vectors and $T$ is the observation horizon. Target construction used branching order zero and one node, so these are deterministic nonlinear extended-path pilots, not stochastic-expectations SEP validation. ``Full-JAX'' recomputes the parameter-dependent steady state and first-order ROM inside the surrogate profile inversion objective. It does not evaluate a full nonlinear SEP likelihood inside HMC. Correction therefore targets the declared surrogate objective, not an exact full-nonlinear posterior. The quick, corrected-resource, and default runs use respectively 1, 2, and 1 chains with 2, 2, and 50 post-warmup draws per chain. These runs establish execution and resource feasibility only.}
\end{table}


The pilot establishes four narrow facts. The parser, parameter-specific
preflight, nonlinear equation solver, residual network, likelihood, and sampler
compose on the full model. The order-zero extended-path targets satisfy a strict
residual check in this small design. The fitted residual improves held-out state
error over the first-order ROM. Finally, the Python and JAX implementations of
the same surrogate profile inversion objective agree at the audited calibration.

The corrected resource run establishes only that endpoint evaluation is
possible under a tiny budget. It retained four draws in total, achieved 100
percent acceptance, and processed 0.0267 draws per second. Those numbers are not
convergence evidence and are not a speed advantage. The acceptance rate instead
suggests that the static-HMC tuning is far from an efficient production regime.

\section{Required confirmatory evidence}

Table~\ref{tab:submission-gates} defines the minimum evidence needed to promote
the working draft to a submission. Counts are protocol defaults, not universal
publication standards; learning curves, uncertainty intervals, and reference
convergence determine whether they are sufficient.

\begin{table}[t]
\centering
\caption{Hard gates for a credible reduced-model methods submission}
\label{tab:submission-gates}
\small
\begin{tabular}{p{0.20\textwidth}p{0.70\textwidth}}
\toprule
Evidence family & Required result \\
\midrule
Scope and provenance & Full 23-equation source verified by content hash and upstream revision; frozen synthetic-data generator and estimand; clean repository revision; complete environment, data, command, seed, and artifact manifests. \\
Nonlinear solution & Positive-order stochastic SEP with multiple nodes and all shocks active; raw attempted and successful counts; scaled residuals; regular, posterior, and boundary-stress strata. \\
Independent parity & Common case identifiers evaluated in Python and MacroModelling.jl; both failures retained in the denominator; QZ, trajectory, and likelihood errors reported on common successful cases. \\
Surrogate accuracy & Disjoint parameter, shock, and trajectory partitions; shrinkage disabled or fitted on a separate calibration partition; untouched-test per-output and tail errors; paired improvement over ROM with a 95 percent interval; learning-curve and out-of-distribution results. \\
Posterior fidelity & Two-parameter full nonlinear grid with resolution-doubling convergence; posterior moments, intervals, and distribution compared with the surrogate and corrected samplers using numerical or Monte Carlo uncertainty. \\
Sampling & Four tuned chains for the scaling experiment; rank-normalized split $\widehat R$, bulk and tail ESS, failures, ESS per gradient evaluation, and ESS per second. \\
Timing & Same-work, same-precision comparison; synchronized warm calls and independent end-to-end processes; compilation, transfers, training, host callbacks, and peak memory reported separately. \\
\bottomrule
\end{tabular}

\vspace{0.4em}
\parbox{0.92\textwidth}{\footnotesize\textit{Notes:} These are claim-specific design requirements. Prespecified thresholds and case lists must be frozen before confirmatory execution. Passing them supports the reduced Gal\'i experiment only; it does not establish universal accuracy across DSGE models.}
\end{table}


All thresholds and case identifiers will be frozen before the final runs. The
evidence builder reads raw result files, verifies the full-model identity, hashes
its sources, and emits both machine-readable and manuscript inputs. A failed
hard gate remains visible in the replication package; it is not converted into
a missing observation or silently omitted path.

\section{Limitations and interpretation}

The preserved evidence is an engineering pilot. Its nonlinear targets use
branching order zero and one node, so it does not validate stochastic
expectations. Most likelihood calls hold the reference steady state and ROM
fixed. The parameter-dependent full-JAX run still targets the surrogate profile
inversion objective, assigns a finite $-10^{12}$ sentinel to failed solves, and
uses a host-callback Schur solve. An exact structural posterior instead requires
explicit support handling and an economically defined likelihood. The sampler
budgets are too small for convergence diagnostics, the data are synthetic and short,
and no independent Julia trajectory or likelihood comparison is yet archived.
The current evidence therefore supports neither exact nonlinear posterior
sampling nor an end-to-end speedup claim.

These limitations determine the paper's next experiments rather than merely
qualifying its conclusion. A credible methods result requires the two-parameter
reference experiment to pass before scaling the posterior dimension. If the
surrogate does not preserve the grid posterior, or if its full accounting does
not improve ESS per second, the negative result should be reported and the claim
reduced to numerical infrastructure.

\section{Conclusion}

Residual learning can be useful for nonlinear DSGE estimation because a
first-order model provides a disciplined baseline and the network concentrates
on the nonlinear discrepancy. Target-corrected HMC can remove error from the
trajectory force, but only relative to the density actually evaluated at the
Metropolis endpoints. The preserved Gal\'i pilot shows that the software path is
feasible and identifies the dominant remaining costs. Publication readiness now
depends on stochastic-SEP validation, an independently converged nonlinear
posterior reference, cross-language parity, efficient multi-chain sampling, and
same-work timing. The accompanying evidence protocol makes those requirements
falsifiable and reproducible.

\section*{Disclosure on computational assistance}

Automated coding and language tools were used to inspect software, generate
tests, and edit the working draft. Every numerical claim in the paper is linked
to a machine-readable source file and must be verified by the authors. The final
submission will follow the target journal's disclosure policy.

\bibliographystyle{plainnat}
\bibliography{references}

\end{document}
